\begin{recipe}{Zelfgemaakte mozzarella}
\label{recipe:mozzarella}
\kicker[Brood,Vegetarisch,Basisrecept]{Brood \& basis · vegetarisch · basisrecept}
\index[register]{Mozzarella}
\index[register]{Melk}
\index[register]{Azijn}

\meta{45 minuten \dvd Circa 250 g mozzarella}

\lettrine{Z}{elf} mozzarella maken klinkt ingewikkelder dan het is: met melk en
azijn en een beetje geduld heb je binnen drie kwartier een verse bol kaas. Het
is ook leuk om samen met kinderen te doen, ze zien meteen hoe wrongel uit melk
ontstaat. Bij het rekken wordt de wei wel behoorlijk heet, dus zorg voor een
schuimspaan en, als je die hebt, vuurvaste handschoenen.

\blockrule{Ingrediënten}
\begin{ingredients}
  \ing{2 liter}{volle melk (gepasteuriseerd)}\index[register]{Melk}
  \ing{7 eetlepels}{witte natuurazijn}\index[register]{Azijn}
  \ingb{zout naar smaak, optioneel}
\end{ingredients}

\blockrule{Bereiding}
\tip{Een digitale keukenthermometer is onmisbaar om de temperaturen van de
melk en de wei precies aan te houden.}
\begin{steps}
  \item Verwarm de melk tot 46\,°C. Haal de pan van het vuur en voeg de
        witte azijn toe. Roer rustig door zodat alles goed mengt. Doe een
        deksel op de pan en laat 10 minuten staan.
  \item Schep de gevormde wrongel met een schuimspaan uit de wei. Knijp
        zoveel mogelijk vocht uit de wrongel en vorm er bolletjes van.
  \item Verwarm de achtergebleven wei tot 65\,°C. Leg één mozzarellabol
        tegelijk 30 seconden in de warme wei op een schuimspaan.
  \item Haal de bol uit de wei en rek en vouw de kaas tot ze glad en
        elastisch wordt. Herhaal het verwarmen en rekken 2 tot 4 keer, tot
        de mozzarella een mooie, gladde structuur heeft.
  \item Vorm de mozzarella tot een mooie bol. Voeg eventueel naar smaak wat
        zout toe.
  \item Leg de mozzarella 10 tot 15 minuten in een kom met ijswater om de
        vorm te behouden en af te koelen.
  \item Bewaar de mozzarella in het water in de koelkast. De kaas blijft
        ongeveer een week goed.
\end{steps}
\end{recipe}
